\begin{proof}[Proof of \cref{obj:lem:copula-frame-legality}]
Fix an admissible tuple and ranks satisfying the hypotheses, and fix an arbitrary realization of all coefficient signs. Admissibility gives
\[
1<p\le2,\qquad
0<\alpha,\beta\le1,\qquad
\frac14\le\gamma\le1.
\]
We first treat the finite-moment construction.

\begin{enumerate}
\item \textit{Tuning and construction inputs.}
The rank assumptions imply \(K\ge32\), so the displayed amplitudes satisfy
\[
0<a\le\frac1{16},\qquad
0<b\le\frac1{256},\qquad
aK^\alpha=\frac1{16},\qquad
bK^\beta=\frac1{256}.
\]
Indeed, \(K^{-\alpha},K^{-\beta}\in(0,1]\), and multiplication by the corresponding positive powers of \(K\) cancels the exponents. Also \(q>0\) and \(0<16b\le1/16<1\), whence
\[
0<\varepsilon=(16b)^{1/q}<1,\qquad
L=\varepsilon^{-1/p}>0.
\]
The tuning identities are
\[
\varepsilon L^p
=\varepsilon\varepsilon^{-1}=1,
\qquad
\varepsilon Lu
=\frac{\varepsilon^{\,1-1/p}}{16}
=\frac{\varepsilon^q}{16}
=b.
\]
Thus the finite-moment tuning meets the input conditions of
\cref{obj:def:copula-frame-priors}. The second tuning also meets those conditions, since
\[
0<2b\le\frac1{128}\le\frac1{16},
\qquad
0<\frac12<1,\qquad 1>0.
\]
For that tuning,
\[
\varepsilon Lu=\frac12\cdot1\cdot2b=b,
\qquad
\varepsilon L^2=\frac12.
\]
% lean: CausalSmith.Stat.FinitepHomogeneityDensegamma.legality_copula_domains

\item \textit{Validity of the conditional table and exact moments.}
Use the coordinates and signs of \cref{obj:def:copula-frame-priors}, and define the two coefficient fields by
\[
\mathcal F_\lambda(x)=\sum_{i=0}^{K}\lambda_i f_i(x),
\qquad
\mathcal F_\eta(x)=\sum_{i=0}^{K}\eta_i f_i(x),
\qquad x\in[0,1].
\]
On a fine cell, write
\[
\theta_i(x)=\frac{\pi(Kx-i)}2,
\qquad
0\le i<K,\quad x\in[i/K,(i+1)/K].
\]
There the only possibly nonzero coordinates are
\(f_i(x)=\cos\theta_i(x)\) and
\(f_{i+1}(x)=\sin\theta_i(x)\). Consequently,
\[
\sum_{i=0}^{K}f_i(x)^2=1,\qquad
|\mathcal F_\lambda(x)|,\ |\mathcal F_\eta(x)|\le2.
\]
The coarse tents have absolute value at most one. Hence their squared-frame interpolation satisfies
\[
|\widetilde g_\sigma(x)|
\le\sum_{i=0}^{K}|g_\sigma(i/K)|f_i(x)^2
\le\sum_{i=0}^{K}f_i(x)^2=1.
\]
For either valid tuning and either branch,
\[
|\xi(x)|\le2a\le\frac18,\qquad
|\upsilon(x)|\le2u\le\frac18,
\qquad
a^2\le\frac1{256},\qquad au\le\frac1{256}.
\]
In the null branch \(t_0=0\). In the alternative branch, positivity of \(1-a^2\) and the preceding bounds give
\[
|t_1(x)|
\le\frac{au\kappa}{1-a^2}
\le\frac1{128},
\]
where the last inequality follows from
\(au\le1/256\), \(\kappa=1/16\), and
\(1-a^2\ge255/256\). Since
\(0\le\xi(x)^2\le1/64\), we also have
\[
|1-\xi(x)^2|\le1,
\qquad
|\zeta(x)|
\le|\xi(x)\upsilon(x)|
   +|t_\nu(x)|\,|1-\xi(x)^2|
\le\frac1{64}+\frac1{128}\le\frac18.
\]
Thus every unmarked and marked bracket in
\cref{obj:def:marked-table} is positive:
\[
1+\ell\xi(x)\ge\frac78,
\qquad
1+\ell\xi(x)+h\upsilon(x)+\ell h\zeta(x)
\ge\frac58,
\qquad \ell,h\in\{-1,1\}.
\]
The frame coordinates are continuous at their cell endpoints, and all three table coordinates are continuous finite expressions in them. The table therefore supplies Borel conditional probability kernels.

Summing the marked probabilities over the outcome sign gives
\[
\sum_{h\in\{-1,1\}}
\frac{\varepsilon}{4}
\{1+\ell\xi(x)+h\upsilon(x)+\ell h\zeta(x)\}
=\frac{\varepsilon}{2}\{1+\ell\xi(x)\}.
\]
Adding the unmarked probability gives
\((1+\ell\xi(x))/2\), and summing over \(\ell\) gives one. In particular, the conditional probability of a mark within either treatment arm is exactly \(\varepsilon\). The marked outcomes have absolute value \(L\), and the unmarked outcome is zero. Therefore
\[
\int_{\mathbb R}|y|^{r_{\mathrm{mom}}}
\,Q_{(1+\ell)/2,P}(dy\mid x)
=\varepsilon L^{r_{\mathrm{mom}}},
\qquad
r_{\mathrm{mom}}\in(0,\infty),\
\ell\in\{-1,1\},\ x\in[0,1].
\]
The tuning identities from the preceding step give the required value \(1\) at order \(p\), and the value \(1/2\) at order two for the bounded tuning.

Summing the signed marked outcomes within each arm also gives
\[
\begin{aligned}
e_P(x)&=\frac{1+\xi(x)}2,\\
m_{0,P}(x)
&=\varepsilon L\frac{\upsilon(x)-\zeta(x)}{1-\xi(x)}
 =\varepsilon L\{\upsilon(x)-t_\nu(x)(1+\xi(x))\},\\
m_{1,P}(x)
&=\varepsilon L\frac{\upsilon(x)+\zeta(x)}{1+\xi(x)}
 =\varepsilon L\{\upsilon(x)+t_\nu(x)(1-\xi(x))\},\\
\tau_P(x)&=2\varepsilon L t_\nu(x).
\end{aligned}
\]
All denominators are positive because \(|\xi(x)|\le1/8\).
% lean: CausalSmith.Stat.FinitepHomogeneityDensegamma.legality_arm_moments

\item \textit{Null membership.}
We record the frame estimates needed for the primitive restrictions. For either coefficient field,
\[
|\mathcal F_\lambda(x)|,\ |\mathcal F_\eta(x)|
\le\sqrt2.
\]
Indeed, on a fine cell its absolute value is at most
\(|\cos\theta_i(x)|+|\sin\theta_i(x)|\), whose square is at most two by
\(\cos^2\theta_i(x)+\sin^2\theta_i(x)=1\) and
\((|\cos\theta_i(x)|-|\sin\theta_i(x)|)^2\ge0\).

For either field, denoted below by
\(\mathcal F\in\{\mathcal F_\lambda,\mathcal F_\eta\}\), the sine and cosine Lipschitz inequalities give, when \(x,z\) belong to the same fine cell,
\[
|\mathcal F(x)-\mathcal F(z)|
\le2|\theta_i(x)-\theta_i(z)|
=\pi K|x-z|.
\]
If \(K|x-z|<1\), the points lie in the same cell or in adjacent cells. In the latter case, split at the common endpoint and add the two cellwise bounds; the two distances sum to \(|x-z|\). Equality of the points gives zero, and reversing their order leaves the bound unchanged. Thus
\[
|\mathcal F(x)-\mathcal F(z)|
\le\pi K|x-z|
\qquad\text{when }K|x-z|<1.
\]
For \(0<s_{\mathrm H}\le1\), the near case uses
\(K|x-z|\le(K|x-z|)^{s_{\mathrm H}}\) and \(\pi<4\); the complementary case uses
\(|\mathcal F(x)-\mathcal F(z)|\le2\sqrt2\le4\) and
\((K|x-z|)^{s_{\mathrm H}}\ge1\). Together these give
\[
|\mathcal F(x)-\mathcal F(z)|
\le4K^{s_{\mathrm H}}|x-z|^{s_{\mathrm H}},
\qquad
0<s_{\mathrm H}\le1,\quad x,z\in[0,1].
\]

The membership calculation applies to valid construction inputs satisfying
\[
\varepsilon Lu=b,\qquad \varepsilon L^p\le10.
\]
In the null branch, the table identities reduce to
\[
e_P(x)=\frac{1+a\mathcal F_\lambda(x)}2,
\qquad
m_{0,P}(x)=b\mathcal F_\eta(x),
\qquad
\tau_P(x)=0.
\]
The frame envelope and amplitude bounds imply
\[
|\xi(x)|\le a\sqrt2\le\frac18,\qquad
\frac14\le e_P(x)\le\frac34,\qquad
|m_{0,P}(x)|\le b\sqrt2\le\frac12.
\]
The frame Hölder estimates, together with the amplitude-power identities, give
\[
\begin{aligned}
|e_P(x)-e_P(z)|
&\le2aK^\alpha|x-z|^\alpha
=\frac18|x-z|^\alpha
\le20|x-z|^\alpha,\\
|m_{0,P}(x)-m_{0,P}(z)|
&\le4bK^\beta|x-z|^\beta
=\frac1{64}|x-z|^\beta
\le20|x-z|^\beta.
\end{aligned}
\]
The effect is continuous, has supremum norm zero, and has zero Hölder differences. These estimates establish
\cref{obj:ass:overlap,obj:ass:propensity-smoothness,obj:ass:baseline-smoothness,obj:ass:effect-smoothness,obj:ass:baseline-cap,obj:ass:effect-cap}.
The normalized conditional table is integrated against \(\lambda\), so
\(P_X=\lambda\), as required by \cref{obj:ass:uniform}. The exact conditional-moment formula yields
\[
\int_{\mathbb R}|y|^p Q_{(1+\ell)/2,P}(dy\mid x)
=\varepsilon L^p\le10,
\]
and hence \cref{obj:ass:raw-moment}. Finally, \(\tau_P=0\) supplies
\cref{obj:ass:null-constancy} with constant zero. Therefore
\(P\in H_0(v)\) by \cref{obj:def:null}. The finite-moment tuning satisfies the two displayed input conditions by the first step.
% lean: CausalSmith.Stat.FinitepHomogeneityDensegamma.legality_null_inNull

\item \textit{Alternative membership.}
Define the positive correction amplitude by
\[
C_{\mathrm{cop}}=\frac{ab\kappa}{1-a^2}>0.
\]
The amplitude bounds give
\[
1-a^2\ge\frac12,\qquad
C_{\mathrm{cop}}\le2ab\kappa\le\frac1{32768}.
\]
To control its smoothness factors, first observe that \(p-1\le1\) and \(q\le1/2\). Consequently,
\[
2\alpha\le\frac{2\alpha}{p-1},
\qquad
2\beta\le\frac{\beta}{q},
\qquad
2S+\frac{S}{2\gamma}\le F(p)<1.
\]
If \(S\ge\gamma\), the left side would be at least
\(2\gamma+1/2\ge1\). Thus
\[
S<\gamma,\qquad \beta\le\gamma.
\]
Rank compatibility and positivity of \(\gamma\) give
\(M^\gamma\le K^S\). Using \(M\ge1\) and the amplitude identities, we obtain
\[
\begin{aligned}
ab&=\frac{K^{-S}}{4096},\\
abM^\gamma&\le\frac{K^{-S}K^S}{4096}=\frac1{4096},\\
abM^\beta&\le abM^\gamma\le\frac1{4096},\\
a^2bK^\beta&=\frac{a^2}{256}\le\frac1{65536}.
\end{aligned}
\]
Multiplication by \(2\kappa=1/8\) now gives
\[
C_{\mathrm{cop}}M^\gamma\le\frac1{32768},
\qquad
C_{\mathrm{cop}}M^\beta\le\frac1{32768},
\qquad
C_{\mathrm{cop}}aK^\beta\le\frac1{524288}.
\]
The frame estimates also give
\[
|\xi(x)|\le\frac18,\qquad
|1+\xi(x)|\le2,\qquad
|\xi(x)-\xi(z)|
\le4aK^\beta|x-z|^\beta.
\]

For valid inputs with \(\varepsilon Lu=b\), the alternative table identities become
\[
m_{0,P}(x)
=b\mathcal F_\eta(x)
 +C_{\mathrm{cop}}\widetilde g_\sigma(x)(1+\xi(x)),
\qquad
\tau_P(x)=-2C_{\mathrm{cop}}\widetilde g_\sigma(x).
\]
Their uniform bounds are
\[
|m_{0,P}(x)|
\le2b+2C_{\mathrm{cop}}\le\frac12,
\qquad
|\tau_P(x)|\le2C_{\mathrm{cop}}\le\frac12.
\]
The design remains \(\lambda\), because the alternative conditional table is normalized. Its propensity is the same function
\((1+a\mathcal F_\lambda)/2\) as in the null branch, so the preceding overlap and propensity-smoothness estimates apply.

For the remaining smoothness estimates, extend the triangular tent by
\[
b_{\triangle,\mathrm{ext}}(r_\triangle)
=\max\{0,\,2\min(r_\triangle,1-r_\triangle)\},
\qquad r_\triangle\in\mathbb R.
\]
This agrees with \(b_\triangle\) on \([0,1]\) and vanishes outside that interval. The elementary Lipschitz inequalities for minimum and maximum give
\[
|b_{\triangle,\mathrm{ext}}(r_\triangle)
 -b_{\triangle,\mathrm{ext}}(s_\triangle)|
\le2|r_\triangle-s_\triangle|,
\qquad r_\triangle,s_\triangle\in\mathbb R.
\]
Define each signed pair contribution and its active index set by
\[
\begin{aligned}
\mathfrak g_j(x)
&=\sigma_j\bigl\{
b_{\triangle,\mathrm{ext}}(Mx-(2j-2))
-b_{\triangle,\mathrm{ext}}(Mx-(2j-1))
\bigr\},\\
\mathcal J_\sigma(x)
&=\{j\in\{1,\ldots,M/2\}:\mathfrak g_j(x)\ne0\},
\qquad x\in[0,1].
\end{aligned}
\]
Thus \(g_\sigma=\sum_{j=1}^{M/2}\mathfrak g_j\).
Each contribution satisfies
\[
|\mathfrak g_j(x)-\mathfrak g_j(z)|
\le4M|x-z|.
\]
The interiors of the pair supports are disjoint, so
\(|\mathcal J_\sigma(x)|\le1\). Only indices in
\(\mathcal J_\sigma(x)\cup\mathcal J_\sigma(z)\) contribute to a difference, and this union has cardinality at most two. Hence
\[
|g_\sigma(x)-g_\sigma(z)|\le8M|x-z|.
\]

On a fine cell, the squared-frame formula is
\[
\widetilde g_\sigma(x)
=g_\sigma(i/K)\cos^2\theta_i(x)
 +g_\sigma((i+1)/K)\sin^2\theta_i(x).
\]
For \(x,z\) in that cell, subtraction and
\(\cos^2+\sin^2=1\) give
\[
\begin{aligned}
\widetilde g_\sigma(x)-\widetilde g_\sigma(z)
={}&\bigl(g_\sigma((i+1)/K)-g_\sigma(i/K)\bigr)\\
&\quad{}\times
\bigl(\sin\theta_i(x)-\sin\theta_i(z)\bigr)
\bigl(\sin\theta_i(x)+\sin\theta_i(z)\bigr).
\end{aligned}
\]
The sine Lipschitz inequality and its unit envelope therefore imply
\[
|\widetilde g_\sigma(x)-\widetilde g_\sigma(z)|
\le\frac{8M}{K}\,\pi K|x-z|
\le32M|x-z|.
\]
Adjacent cell formulas agree at their common endpoint. Splitting any ordered interval at the intervening fine-grid endpoints and summing the cellwise estimates proves this bound for all \(x,z\in[0,1]\); equal points give zero, and reversing their order preserves the bound.

If \(M|x-z|\le1\), use
\(M|x-z|\le(M|x-z|)^{s_{\mathrm H}}\). If \(M|x-z|>1\), use the envelope
\(|\widetilde g_\sigma(x)-\widetilde g_\sigma(z)|\le2\). Thus
\[
|\widetilde g_\sigma(x)-\widetilde g_\sigma(z)|
\le32M^{s_{\mathrm H}}|x-z|^{s_{\mathrm H}},
\qquad 0<s_{\mathrm H}\le1.
\]
The same inequality holds for \(s_{\mathrm H}=0\), with the convention
\(r^0=1\) for \(r\ge0\), since the left side is at most two.

The baseline product difference decomposes as
\[
\begin{aligned}
m_{0,P}(x)-m_{0,P}(z)
={}&b\{\mathcal F_\eta(x)-\mathcal F_\eta(z)\}\\
&+C_{\mathrm{cop}}\bigl[
 \{\widetilde g_\sigma(x)-\widetilde g_\sigma(z)\}(1+\xi(x))
 +\widetilde g_\sigma(z)\{\xi(x)-\xi(z)\}
\bigr].
\end{aligned}
\]
Applying the displayed estimates to these three terms gives
\[
\begin{aligned}
|m_{0,P}(x)-m_{0,P}(z)|
&\le
\bigl(4bK^\beta
 +64C_{\mathrm{cop}}M^\beta
 +4C_{\mathrm{cop}}aK^\beta\bigr)|x-z|^\beta\\
&\le
\left(\frac4{256}+\frac{64}{32768}
                  +\frac4{524288}\right)|x-z|^\beta\\
&\le20|x-z|^\beta.
\end{aligned}
\]
Similarly,
\[
\begin{aligned}
|\tau_P(x)-\tau_P(z)|
&\le64C_{\mathrm{cop}}M^\gamma|x-z|^\gamma\\
&\le\frac{64}{32768}|x-z|^\gamma
\le20|x-z|^\gamma.
\end{aligned}
\]
All these representatives are continuous. Their uniform bounds and Hölder estimates establish the baseline and effect restrictions in
\cref{obj:ass:baseline-smoothness,obj:ass:effect-smoothness,obj:ass:baseline-cap,obj:ass:effect-cap}.
The conditional-moment calculation supplies
\cref{obj:ass:raw-moment} whenever \(\varepsilon L^p\le10\).
Together with the design, overlap, and propensity restrictions already established, this proves
\(P\in\mathcal M_v\) by \cref{obj:def:model}. In particular it applies to the finite-moment tuning.
% lean: CausalSmith.Stat.FinitepHomogeneityDensegamma.legality_alternative_inModel

\item \textit{Centering of the interpolated tents.}
Choose the dyadic exponents and the refinement factor so that
\[
K=2^{k_{\mathrm{rank}}},\qquad
M=2^{m_{\mathrm{rank}}},
\qquad
k_{\mathrm{rank}},m_{\mathrm{rank}}\in\mathbb N,
\qquad
1\le m_{\mathrm{rank}}\le k_{\mathrm{rank}},
\]
\[
r_{\mathrm{grid}}=2^{k_{\mathrm{rank}}-m_{\mathrm{rank}}}
=\frac KM\in\mathbb N,\qquad r_{\mathrm{grid}}\ge16.
\]
The inequality between the dyadic exponents follows from \(M\le K\), and \(M\ge2\) implies that \(M\) is even.

Every coarse cell has the same fine-grid sample sum of its unsigned tent. More precisely, support of the extended tent and the change of index
\(i=i_{\mathrm c}r_{\mathrm{grid}}+h_{\mathrm{grid}}\) give
\[
\sum_{i=0}^{K}
 b_{\triangle,\mathrm{ext}}(Mi/K-i_{\mathrm c})
=
\sum_{h_{\mathrm{grid}}=0}^{r_{\mathrm{grid}}}
 b_\triangle(h_{\mathrm{grid}}/r_{\mathrm{grid}}),
\qquad
i_{\mathrm c}\in\{0,\ldots,M-1\}.
\]
The two cells in each signed pair therefore have equal sample sums and opposite signs. Consequently,
\[
\sum_{i=0}^{K}g_\sigma(i/K)=0,
\qquad g_\sigma(0)=g_\sigma(1)=0.
\]
In particular,
\[
\sum_{i=0}^{K-1}g_\sigma(i/K)=0,
\qquad
\sum_{i=0}^{K-1}g_\sigma((i+1)/K)=0.
\]

Integrating the sine and cosine squares over one fine cell, by an affine change of variable, gives
\[
\int_{i/K}^{(i+1)/K}\cos^2\theta_i(x)\,dx
=
\int_{i/K}^{(i+1)/K}\sin^2\theta_i(x)\,dx
=\frac1{2K}.
\]
The cell formula therefore yields
\[
\begin{aligned}
\int_0^1\widetilde g_\sigma(x)\,\lambda(dx)
&=\sum_{i=0}^{K-1}
 \frac{g_\sigma(i/K)+g_\sigma((i+1)/K)}{2K}\\
&=0.
\end{aligned}
\]
% lean: CausalSmith.Stat.FinitepHomogeneityDensegamma.smoothedTent_integral_zero_of_legalityRanks

\item \textit{Second-moment bounds.}
We first sharpen the coarse-field difference estimate used for interpolation error. Define the midpoint between the two cells of a pair by
\[
x_{j,\mathrm{mid}}=\frac{2j-1}{M},
\qquad j\in\{1,\ldots,M/2\}.
\]
On either side of this point, one of the two extended tents in
\(\mathfrak g_j\) vanishes. The tent Lipschitz estimate thus gives a constant \(2M\) when both arguments are on the same side. If
\(x\le x_{j,\mathrm{mid}}\le z\), both tents vanish at the midpoint, and
\[
\begin{aligned}
|\mathfrak g_j(x)-\mathfrak g_j(z)|
&\le|\mathfrak g_j(x)|+|\mathfrak g_j(z)|\\
&\le2M(x_{j,\mathrm{mid}}-x)
     +2M(z-x_{j,\mathrm{mid}})
=2M|x-z|.
\end{aligned}
\]
Reversing the arguments covers the other crossing case. Summing over the at most two active pair indices gives
\[
|g_\sigma(x)-g_\sigma(z)|\le4M|x-z|,
\qquad x,z\in[0,1].
\]

The partition identity gives the exact interpolation-error expression
\[
\widetilde g_\sigma(x)-g_\sigma(x)
=
\sum_{i=0}^{K}
 \{g_\sigma(i/K)-g_\sigma(x)\}f_i(x)^2.
\]
Whenever \(f_i(x)\ne0\), its adjacent-cell support implies
\(|i/K-x|\le1/K\); terms with \(f_i(x)=0\) vanish. Hence
\[
\begin{aligned}
|\widetilde g_\sigma(x)-g_\sigma(x)|
&\le\sum_{i=0}^{K}\frac{4M}{K}f_i(x)^2\\
&=\frac{4M}{K}\le\frac14.
\end{aligned}
\]

For the exact coarse second moment, direct integration gives
\[
\int_0^1 b_\triangle(r_\triangle)^2\,dr_\triangle
=2\int_0^{1/2}(2r_\triangle)^2\,dr_\triangle
=\frac13.
\]
Each coarse cell therefore contributes \(1/(3M)\). The cell interiors are disjoint, the signs square to one, and the even rank supplies all \(M\) cells, so
\[
\int_0^1g_\sigma(x)^2\,\lambda(dx)=\frac13.
\]
Pointwise, the error estimate and the elementary square inequality give
\[
\begin{aligned}
g_\sigma(x)^2
&\le2\widetilde g_\sigma(x)^2
   +2\{\widetilde g_\sigma(x)-g_\sigma(x)\}^2\\
&\le2\widetilde g_\sigma(x)^2+\frac18.
\end{aligned}
\]
Integration against the probability measure \(\lambda\) yields
\[
\frac13
\le2\int_0^1\widetilde g_\sigma(x)^2\,\lambda(dx)+\frac18,
\]
which implies the lower bound below. The upper bound follows from the unit envelope:
\[
\frac1{16}
\le\int_0^1\widetilde g_\sigma(x)^2\,\lambda(dx)
\le\int_0^1 1\,\lambda(dx)=1.
\]
% lean: CausalSmith.Stat.FinitepHomogeneityDensegamma.smoothedTent_second_moment_lower_of_legalityRanks

\item \textit{Effect identity and distance calibration.}
Define the effect multiplier by
\[
A_{\mathrm{eff}}
=2C_{\mathrm{cop}}
=\frac{2ab\kappa}{1-a^2}>0.
\]
The alternative table calculation gives, pointwise,
\[
\tau_P(x)=-A_{\mathrm{eff}}\widetilde g_\sigma(x)
=-\frac{2ab\kappa}{1-a^2}\widetilde g_\sigma(x).
\]
Since the interpolated tent is centered,
\[
\int_0^1\tau_P(x)\,\lambda(dx)
=-A_{\mathrm{eff}}
 \int_0^1\widetilde g_\sigma(x)\,\lambda(dx)=0.
\]
Substitution into the definition of \(d(P)\), followed by factoring the nonnegative multiplier out of the square root, gives
\[
d(P)
=A_{\mathrm{eff}}
 \left(\int_0^1\widetilde g_\sigma(x)^2\,\lambda(dx)\right)^{1/2}.
\]

The amplitude product and \(\kappa=1/16\) give
\[
2ab\kappa=2^{-15}K^{-S}.
\]
Because \(0\le a^2\le1/256\), the denominator is positive and lies between \(1/2\) and one. Therefore
\[
2^{-15}K^{-S}
\le A_{\mathrm{eff}}
\le2^{-14}K^{-S}.
\]
Taking square roots of the second-moment bounds gives
\[
\frac14
\le
\left(\int_0^1\widetilde g_\sigma(x)^2\,\lambda(dx)\right)^{1/2}
\le1.
\]
Multiplication of the two intervals proves
\[
2^{-17}K^{-S}\le d(P)\le2^{-14}K^{-S}.
\]

For the strict comparison, \cref{obj:lem:causal-nonempty} supplies
\(d_0=1/(8\sqrt2)\). Since \(S>0\) and \(K\ge1\), we have
\(K^{-S}\le1\). Also \(\sqrt2<2\), so
\[
d(P)\le2^{-14}K^{-S}
\le\frac1{16384}
<\frac1{8\sqrt2}=d_0.
\]
This establishes the effect, centering, and distance conclusions for every alternative realization in the finite-moment construction.
% lean: CausalSmith.Stat.FinitepHomogeneityDensegamma.legality_separation_of_moments

\item \textit{Bounded construction.}
Now take an admissible \(w\) and \(v=(2,\alpha,\beta,\gamma)\). This \(v\) is admissible, and the assumed phase and rank conditions are exactly those used above. Under
\[
u=2b,\qquad \varepsilon=\frac12,\qquad L=1,
\]
the input verification in the first step applies, and
\[
\varepsilon Lu=b,\qquad
\varepsilon L^p=\varepsilon L^2=\frac12\le10.
\]
Thus the null-membership calculation gives all restrictions of
\(H_0(v)\), and the alternative-membership calculation gives all restrictions of \(\mathcal M_v\).

In either branch, the normalized table has outcome support
\[
Y\in\{-1,0,1\},
\qquad P(|Y|\le1)=1.
\]
Indeed, each conditional table atom has absolute outcome at most one, so its conditional probability of this event is one; integrating over \(\lambda\) preserves that value. This supplies
\cref{obj:ass:bounded-outcome}. Together with the previously established restrictions, it gives bounded null and alternative membership by
\cref{obj:def:bounded-null,obj:def:bounded-model}.

The table effect calculation uses \(\varepsilon Lu=b\), so every bounded alternative has the same displayed effect formula. The tent centering and second-moment estimates depend on the ranks and coarse signs, and the distance calibration uses the retained amplitudes. Hence the centering identity, both distance bounds, and the strict comparison with \(d_0\) also apply. Finally, the exact arm-moment calculation gives
\[
\int_{\mathbb R}|y|^2
\,Q_{(1+\ell)/2,P}(dy\mid x)
=\varepsilon L^2=\frac12,
\qquad \ell\in\{-1,1\},\quad x\in[0,1].
\]
The coefficient realization was arbitrary throughout, so all conclusions hold for every realization in both priors.
% lean: CausalSmith.Stat.FinitepHomogeneityDensegamma.copula_frame_legality
\end{enumerate}
\end{proof}
